\chapter*{Preface}
\addcontentsline{toc}{chapter}{Preface}

Modern medicine relies on an extraordinary range of imaging tools---from split-second X-rays and bedside ultrasound to high-field MRI, molecular PET/CT, and microscopic optical eye scans. Yet for most patients and students, the alphabet soup of acronyms (CT, CTA, CTP, DWI, DTI, MRA, MRCP, SPECT, OCTA) can feel overwhelming, especially when a scan has just been ordered for you or a loved one.

This illustrated guide was written to bridge the gap between technical radiology textbooks and everyday patient understanding. Each chapter pairs a visual overview with clear, plain-English explanations of how the technology works, what symptoms or clinical questions prompt a doctor to order it, how to prepare, what the exam feels like, and how specialists interpret normal and abnormal patterns.

\section*{How to Use This Book}
\begin{itemize}
  \item \textbf{Organized by Physical Energy (Parts I--VII):} After an introductory classification chapter, the book groups modalities into seven parts: X-Ray, Magnetic Field, Ultrasound, Nuclear Medicine, Optical, Photoacoustic and Electromagnetic, and Hybrid Imaging.
  \item \textbf{Start With Your Symptoms or Body Area:} If you are unsure which chapter applies to your situation, turn to \textbf{Appendix~\ref{app:quick-finder}} (\textit{Quick-Finder Guide}) at the back of the book to jump directly from common symptoms to the relevant chapters.
  \item \textbf{Compare Radiation and Safety:} Turn to \textbf{Appendix~\ref{app:radiation-chart}} (\textit{Radiation Dose and Safety Comparison Chart}) for a side-by-side comparison of zero-radiation modalities and typical ionizing-radiation doses expressed in everyday natural-background equivalents.
  \item \textbf{Clinical vs.\ Research-Stage Modalities:} Everyday hospital staples (such as X-ray, CT, MRI, ultrasound, PET/CT, and OCT) include full patient-visit sections covering preparation, what happens during the scan, and follow-up. Family overview chapters and strictly investigational research technologies (such as magnetic particle imaging, terahertz imaging, microwave imaging, and Raman spectroscopic imaging) are clearly distinguished so readers always know which tools are proven clinical standards and which remain in laboratory or prototype testing.
  \item \textbf{References and Further Reading:} Every chapter separates peer-reviewed clinical guidelines and institutional medical references (\textit{Key References}) from accessible general encyclopedic overviews (\textit{Further Reading}).
\end{itemize}
